% ---------------------------------------------------------------
\section{Hollow Moon}\label{sec:hollowmoon}

Everybody argues about the wrong sentence. The popular version of this section is a shouting match
over whether the Moon ``rang like a bell,'' as though that were still open. It is not open. It was
measured. Apollo~12 dropped its ascent stage onto the surface on 20 November 1969 at roughly 1.7
km/s, and the seismometer the crew had left behind kept sending signal for something close to an
hour. Apollo~13's spent S-IVB booster hit in April 1970 with about eleven times that energy and the
reverberation ran longer still. Later impacts and natural events gave the same profile: a slow rise
to peak, then a decay so gentle the trace looks like a mistake to anyone raised on terrestrial
records. Nobody in the field disputes the data. The papers are published, the tapes exist, the
numbers have been re-reduced repeatedly since.

\actualq{Did the Moon ring for an hour?}{Will the people who own the instruments accept an answer
the arithmetic forces on them, when accepting it costs them the model they built their careers
inside?}

\figwrap[r]{0.42}{The Apollo~12 record of the ascent-stage impact --- the trace that will not stop.
The argument was never about whether this happened. It is about what a body has to be made of to
produce it.}{https://commons.wikimedia.org/wiki/File:Seismometer_reading_from_impact_made_by_Lunar_Module_ascent_stage_(S69-59547).jpg}{fig:lunarseismo}

\subsection{What the Damping Number Actually Says}\label{sec:lunardamping}

Strip the mysticism and you are left with one quantity: $Q$, the seismic quality factor, which is
just a measure of how much energy a material keeps per oscillation instead of turning into heat.
Earth's crust runs $Q$ of a few hundred at best; saturated sediment can drop below a hundred, which
is why a quarry blast in Ohio is over in seconds. The Moon returns $Q$ between roughly 3,\!000 and
5,\!000, and in the deep interior some estimates run higher. That is not a small discrepancy. It is
an order of magnitude and change, and it is the whole of the anomaly. \hi{A ball of ordinary rock,
assembled the way we say the Moon was assembled, should not hold a note for an hour.}

The standard answer is that it is dry and broken. No water, no pore fluid, therefore almost nothing
to convert vibration into heat; and a megaregolith kilometres deep, shattered by four billion years
of bombardment, that scatters the wavefront into a long incoherent tail instead of letting it pass
through cleanly. This is a real mechanism. It is also, and this is the part nobody says out loud in
the press release, a model with free parameters --- scattering layer thickness, velocity gradient,
intrinsic versus scattering attenuation --- tuned after the fact until the synthetic seismogram
resembles the one on the tape. It reproduces the coda. It has never independently predicted one.
And the split between how much of the decay is scattering and how much is intrinsic loss is still
being argued in the literature fifty-five years later, which is a strange thing to call settled.

So the honest position is narrower and far more uncomfortable than either camp wants. The cartoon
hollow Moon --- a thin empty shell, a metal balloon --- is dead, and it deserves to be. GRAIL mapped
the gravity field to kilometre resolution in 2012 and it is the field of a differentiated rocky
body: crust of roughly 34--43 km, mantle, a partially molten layer, a small core near 240--350 km.
Apollo's network caught shear waves crossing the interior, and shear waves do not travel through a
vacuum. Laser ranging pins the moment of inertia, and therefore the mass distribution, tighter than
any conspiracy has room to hide in. Fine. Concede all of it. \emph{The damping is still there.} You
have not explained the anomaly by proving the crude version of it wrong; you have only made the
remaining question harder, because now whatever is producing an hour-long coda has to do it inside a
body whose bulk properties you have already measured.

\subsection{The Reason Nobody Will Touch the Math}\label{sec:whynobodytouches}

Here is my claim, and I want it stated plainly so it can be attacked. The reluctance in this field
is not evidential. The data has been sitting in the archive since Nixon. The reluctance is about
what follows if the structural reading turns out to be right, because \emph{artificial} is not a
statement about geology. It is a statement about intent.

A natural Moon is an accident --- a splash of debris from the Theia impact that happened to settle
into orbit, happened to end up tidally locked with one face permanently toward us, happened to
subtend almost exactly the same angular diameter as the Sun so that a species evolving on the
surface below would be handed total eclipses as a gift. Accidents need no author, and an accident
can be published without disturbing anything.

An engineered Moon is a different sentence entirely. It means the tidal lock is a design choice, and
a design choice means one hemisphere is aimed at this planet permanently and deliberately. It means
the thing that stabilised our axial tilt, and therefore our seasons, and therefore the climate
window in which complex life became possible, was placed. It means the Earth is not a rock that
happened to grow biology. It means the Earth is an experiment with an observation deck, and that
someone has been on station over it, unblinking, on the same face, for as long as there has been
anything here worth observing. Twenty-four hours a day. No shift change. That is what the word
\emph{artificial} costs, and it is why the arithmetic gets called fringe instead of getting funded.

\begin{funfact}
Mikhail Vasin and Alexander Shcherbakov published the ``Spaceship Moon'' hypothesis in 1970 in
\emph{Sputnik}, a Soviet popular-science magazine, months after the Apollo~12 coda was recorded.
Their mechanism was wrong and their evidence was thin. But note the timing --- two Soviet writers
looked at an American seismogram and immediately understood that the interesting number was the
decay rate, not the amplitude. That was more than half a century ago, and the decay rate is still
the interesting number.
\end{funfact}

\subsection{Did They Ever Actually Leave?}\label{sec:didtheyleave}

So I will put the question to you the way it deserves to be put, and you can sit with it, because I
am not going to pretend I can close it.

Every culture that left us writing put its gods in the sky and then told us they departed. They
ascended, they withdrew, they went behind the veil, they promised to return. We read those texts as
farewell notes. But a farewell is what it looks like from the ground when something stops introducing
itself. It is not the same as absence. If there is a platform in permanent orbit with one face fixed
upon us --- and the object over your head does exactly that, whatever built it --- then the gods of
old never needed to leave. They only needed to stop talking.

And if that is true, the strange thing is not that the Moon rang. The strange thing is how badly we
want it to have been nothing, and how much professional machinery exists to make sure the question
never gets asked in a room with a budget.
